\scanpage{148}
% archive foliation: 147
% Typescript is upside down in the supplied scan; transcribed in its reading orientation.
\begin{center}- 167 -\end{center}
\hfill n\textsuperscript{o}329

problèmes d'application universelle (\uline{loc.cit.}). Reste à prouver l'existence
du couple $(i,\widehat{X})$.

1) \uline{Définition de} $\widehat{X}$. \uline{Soit} $\widehat{X}$
\uline{l'ensemble des filtres de Cauchy minimaux} (n\textsuperscript{o}2)
sur $X$. Nous allons définir sur $\widehat{X}$ une structure uniforme.
Pour cela, pour tout entourage \uline{symétrique} $V$ de $X$, désignons par
$\widetilde{V}$ \uline{l'ensemble des couples}
$(\underline{X},\underline{Y})$ \uline{de filtres de Cauchy minimaux ayant en commun
un ensemble petit d'ordre} $V$ ; montrons que les ensembles $\widetilde{V}$
forment un système fondamental d'entourages d'une structure uniforme sur
$\widehat{X}$. En effet :

1\textsuperscript{o} Comme tout $\underline{X}\in\widehat{X}$ est un filtre de
Cauchy, on a par définition $(\underline{X},\underline{X})\in\widetilde{V}$
pour tout entourage symétrique $V$ de $X$, donc $(U'_{\mathrm{I}})$ est vérifié.

2\textsuperscript{o} Si $V,V'$ sont deux entourages symétriques de $X$,
$W=V\cap V'$ est un entourage symétrique, et tout ensemble petit d'ordre $W$
est petit d'ordre $V$ et d'ordre $V'$ ; donc on a
$\widetilde{W}\subset\widetilde{V}\cap\widetilde{V}'$, ce qui prouve $(B_{\mathrm{I}})$.

3\textsuperscript{o} Les ensembles $\widetilde{V}$ sont symétriques par définition,
donc $(U'_{\mathrm{II}})$ est vérifié.

4\textsuperscript{o} Etant donné un entourage symétrique $V$ de $W$, soit $W$ un
entourage symétrique tel que $W^2\subset V$. Considérons trois filtres de Cauchy
minimaux $\underline{X},\underline{Y},\underline{Z}$ tels que
$(\underline{X},\underline{Y})\in\widetilde{W}$ et
$(\underline{Y},\underline{Z})\in\widetilde{W}$ ; il existe donc deux ensembles
$M,N$, petits d'ordre $W$ et tels que $M\in\underline{X}\cap\underline{Y}$,
$N\in\underline{Y}\cap\underline{Z}$. Comme $M$ et $N$ appartiennent à
$\underline{Y}$, $M\cap N$ n'est pas vide, donc (n\textsuperscript{o}1,prop.1),
$M\cup N$ est petit d'ordre $W^2\subset V$ ; comme $M\cup N$ appartient à
$\underline{X}$ et à $\underline{Z}$, on a $W^2\subset\widetilde{V}$ ;
d'où $(U'_{\mathrm{III}})$.

Montrons en outre l'espace uniforme $\widehat{X}$ est \uline{séparé}.
En effet, soient $\underline{X},\underline{Y}$ deux filtres de Cauchy minimaux
sur $X$ tels que $(\underline{X},\underline{Y})\in\widetilde{V}$ pour
\uline{tout} entourage symétrique $V$ de $X$. Il est immédiat que les ensembles
$M\cup N$, où $M\in\underline{X}$, $N\in\underline{Y}$ forment la base d'un
filtre $\underline{Z}$ moins fin que $\underline{X}$ et que $\underline{Y}$.
Or, $\underline{Z}$ est \uline{un filtre de Cauchy}, car pour tout entourage
symétrique $V$ de $X$, il y a par hypothèse un ensemble $P$ petit d'ordre $V$
et appartenant à la fois à $\underline{X}$ et à $\underline{Y}$, donc
